\documentclass[12pt,letterpaper]{article}
\usepackage[margin=1in]{geometry}
\usepackage{fontspec}
\setmainfont{Times New Roman}
\setlength{\emergencystretch}{2em}
\usepackage{amsmath,amssymb,hyperref}
\setcounter{secnumdepth}{0}
\title{Supplementary derivations: How Scoring Responses Change Attribution Evaluation in a Diffusion Policy}
\author{}
\date{}
\begin{document}
\maketitle
These derivations preserve the complete analytic arguments accompanying the revised manuscript. They use the same definitions of Q, N, the squared residual r and the all-baseline residual B. Numerical examples here are analytic constructions, not observed robot-policy ranking reversals.
\appendix
\section{Response identities and target-induced reversal}
\label{sec:proofs}

\subsection{Normalized response difference}
\label{sec:response-proof}

Assume exact self-reference and $B>0$. Set $\eta=\sqrt{\epsilon}$, $b=\sqrt{B+\epsilon}$ and $z=\sqrt{r+\epsilon}$, so $B=(b-\eta)(b+\eta)$. Then
\begin{equation}
C_Q-C_N=\frac{b^2-z^2}{B}-\frac{b-z}{b-\eta}
=\frac{(z-\eta)(b-z)}{B}.
\end{equation}
For $\eta\le z\le b$, the product is nonnegative and maximized at $z=(\eta+b)/2$, giving
\begin{equation}
0\le C_Q-C_N\le\frac{(b-\eta)^2}{4B}
=\frac{b-\eta}{4(b+\eta)}\le\frac14.
\end{equation}
For $z>b$, the difference is negative, so the $1/4$ upper bound holds for all $r\ge0$. Nonnegative AUC weights sum to one, yielding $\Delta\le1/4$ for every defined curve. On the nine-point schedule they are $(0.005,0.025,0.045,0.075,0.10,0.15,0.225,0.25,0.125)$. The endpoints contribute zero difference and have total weight 0.13, giving $\Delta\le0.87/4=0.2175$. For positive $\epsilon$ and finite $B$, the fixed-$B$ schedule bound $0.87(b-\eta)/[4(b+\eta)]$ is strictly smaller. Fully in-range curves additionally satisfy $\Delta\ge0$, strictly if a point with positive AUC weight has $0<r<B$. If $B=0$, normalization is undefined. When the 1\% and 5\% language values coincide, dropping the 1\% node changes their combined AUC weight by 0.02; the $1/4$ pointwise bound therefore bounds the paired change by $0.02/4=0.005$ for fully in-range calls.

\subsection{Finite-stabilizer rank reversal}
\label{sec:reversal-proof}

Let $g(x_1,x_2)=x_1^2+\beta x_2$, with $\beta=4/5$, $x'=(0,0)$ and $x=(1,1)$. Along the diagonal, the scalar action residual and its square are
\[
\rho(\alpha)=(1-\alpha)(1+\alpha+\beta),\qquad r(\alpha,\alpha)=\rho(\alpha)^2.
\]
Since $\nabla g=(2\alpha,\beta)$ there, differentiating the targets for $d_a=1$ gives
\begin{equation}
\nabla Q=\rho(\alpha)(2\alpha,\beta),\qquad
\nabla N=\frac{\rho(\alpha)}{\sqrt{\rho(\alpha)^2+\epsilon}}(2\alpha,\beta).
\end{equation}
Polynomial integration yields
\begin{equation}
A^Q=\left(\frac12+\frac\beta3,\frac{2\beta}3+\frac{\beta^2}2\right)
=\left(\frac{23}{30},\frac{64}{75}\right),
\end{equation}
so feature 2 leads by $13/150$. Norm attributions tend to $(1,\beta)$ as $\epsilon$ decreases to zero.

For the positive stabilizer, define $\psi(\rho)=1-\rho/\sqrt{\rho^2+\epsilon}\ge0$. Its derivative is $-\epsilon/(\rho^2+\epsilon)^{3/2}<0$, and $\rho(\alpha)\ge\beta(1-\alpha)$. Therefore
\begin{equation}
\begin{aligned}
\int_0^1\psi(\rho(\alpha))\,d\alpha
&\le\int_0^1\left[1-\frac{\beta(1-\alpha)}{\sqrt{\beta^2(1-\alpha)^2+\epsilon}}\right]d\alpha\\
&=1-\frac{\sqrt{\beta^2+\epsilon}-\sqrt{\epsilon}}{\beta}
\le\frac{\sqrt{\epsilon}}{\beta}.
\end{aligned}
\end{equation}
The identities $1-A_1^N=\int_0^1 2\alpha\psi(\rho(\alpha))\,d\alpha$ and $\beta-A_2^N=\beta\int_0^1\psi(\rho(\alpha))\,d\alpha$, together with $0\le2\alpha\le2$ and $\epsilon=10^{-12}$, give
\begin{equation}
0\le1-A_1^N\le2.5\times10^{-6},\qquad
0\le\beta-A_2^N\le10^{-6}.
\end{equation}
Thus $A_1^N-A_2^N\ge0.1999975$. Every coordinate is positive, so ordinary and absolute-coordinate rankings agree. Only the target changes. At the half-feature insertion point, feature 1 first gives normalized Q response about 0.80 versus 0.69 for feature 2 first, and N response about 0.56 versus 0.44. Deletion swaps these intermediate values. Hence feature 1 first scores better under both responses, despite exact Q-target IG ranking feature 2 first.

\subsection{Finite-grid random-order calibration}
\label{sec:random-calibration}

For $g(x)=\sum_{i=1}^n x_i$, an all-one input and zero baseline, every feature order gives residual distances $k$ after $k$ deletions and $n-k$ after $k$ insertions. At $q=k/n$, normalized Q curves are $1-q^2$ and $1-(1-q)^2$. Both have second derivative $-2$, so their trapezoidal areas underestimate their common integral $2/3$ by $1/(6n^2)$ on the full-prefix grid. Their identical curvature also gives identical quadrature errors on any common grid, including an asymmetric one. The nine-point schedule with $n=100$ gives area 0.65976. Unstabilized N gives the lines $1-q$ and $q$, both with trapezoidal area 0.5 on either schedule.

\subsection{A common response change can reorder rankings}
\label{sec:ranking-comparison}

Take three features and insertion fractions $\{0,1/3,2/3,1\}$, with $B=1$. Let orders $(1,2,3)$ and $(3,2,1)$ give residual sequences $(1,1/4,1/4,0)$ and $(1,16/25,0,0)$, respectively. Their intermediate subsets are distinct, so these values can belong to one scalar policy with the same endpoints. Q insertion AUCs are $2/3$ and $31/50=0.620$; with $\epsilon=0$, N AUCs are $1/2$ and $17/30\approx0.567$. Thus the better order reverses. The inequalities persist at $\epsilon=10^{-12}$. This establishes possibility, not a measured comparison of the saved cohorts.

\section{A smooth counterexample to integration-grid agreement}
\label{sec:aliasing}

\subsection{Unresolved coordinate variation}
\label{sec:aliasing-proof}

For an integer $K>0$, define
\begin{align}
h(\alpha)&=\frac25\sin^2(\pi K\alpha),\\
g(x_1,x_2)&=\frac25x_1+\frac35x_2+
(x_1-x_2)h\!\left(\frac{x_1+x_2}{2}\right),\\
Q(x_1,x_2)&=-\frac12\bigl(g(x_1,x_2)-1\bigr)^2.
\end{align}
Take baseline $(0,0)$, actual input $(1,1)$, and fixed reference $g(1,1)=1$.
The transverse term vanishes along $(\alpha,\alpha)$, so
$g(\alpha,\alpha)=\alpha$ and $Q(\alpha,\alpha)=-(1-\alpha)^2/2$.
Differentiating before restricting to that path gives
\begin{equation}
\nabla Q(\alpha,\alpha)=(1-\alpha)
\left(\frac25+h(\alpha),\frac35-h(\alpha)\right).
\label{eq:aliasing-gradient}
\end{equation}
Thus the oscillation transfers gradient between coordinates while leaving their
sum unchanged.

Use the composite trapezoidal rule with $m$ intervals.
Whenever $m$ divides $K$, every node $\alpha=j/m$ has $h(\alpha)=0$.
The sampled coordinate functions are then linear. Both the trapezoidal rule and endpoint-inclusive average integrate a linear function exactly on any uniform grid, giving
\begin{equation}
A_m^{\mathrm{trap}}=\left(\frac15,\frac3{10}\right),\qquad
\mathbf{1}^{\mathsf T}A_m^{\mathrm{trap}}=\frac12=Q(1,1)-Q(0,0).
\label{eq:aliasing-coarse}
\end{equation}
Here $\mathbf{1}=(1,1)^{\mathsf T}$ sums the coordinates. Every such grid has zero
completeness residual and ranks feature 2 first, giving identical masks and
perturbation curves under a deterministic ranking rule.

The exact coordinate integrals differ. Since $h(1-\alpha)=h(\alpha)$,
substitution gives
\[
\int_0^1\alpha h(\alpha)\,d\alpha
=\int_0^1(1-\alpha)h(\alpha)\,d\alpha.
\]
Adding these expressions and using $\int_0^1 h(\alpha)\,d\alpha=1/5$
shows that either integral is $1/10$. Consequently,
\begin{equation}
A_{\mathrm{exact}}^Q=\left(\frac3{10},\frac15\right).
\label{eq:aliasing-exact}
\end{equation}
The exact order puts feature 1 first. Coarse errors $(-1/10,1/10)$ cancel in
completeness, but their $L_1$ norm, $1/5$, is 40\% of the exact attribution's norm.

On the $2K$ grid, $h$ alternates between zero and $2/5$, with trapezoidal integral
$1/5$. Reflection symmetry gives area $1/10$ for $(1-\alpha)h$, recovering
Equation~\ref{eq:aliasing-exact}. This recovery is construction-specific. For
finitely many uniform budgets, a common multiple $K$ hides every oscillation.
Agreement cannot certify coordinate integrals without a regularity bound on unresolved variation.

The endpoint-inclusive average has the same coarse failure.
When $m$ divides $K$, the average of the sampled linear functions also yields
Equation~\ref{eq:aliasing-coarse}. On the $2K$ grid, only odd nodes contribute
to $(1-\alpha)h$. Their sum is
$\frac25\sum_{\ell=0}^{K-1}(1-(2\ell+1)/(2K))=K/5$.
Dividing by the $2K+1$ nodes gives
\begin{equation}
A_{2K}^{\mathrm{avg}}=
\left(\frac15+\frac{K}{5(2K+1)},\
\frac3{10}-\frac{K}{5(2K+1)}\right).
\label{eq:aliasing-average}
\end{equation}
The first coordinate exceeds the second for $K>1/2$, including every allowed
integer. Coordinates remain inexact; exact $2K$ recovery belongs to the trapezoidal rule.

\subsection{Exact IG and perturbation quality}
\label{sec:aliasing-quality}

Set $K=128$, so $h(1/2)=0$ and the whole-feature interventions give
$g(1,0)=2/5$ and $g(0,1)=3/5$. Under the same normalized Q response
and fractions $\{0,1/2,1\}$, feature-2-first insertion values are
$(0,0.84,1)$ and deletion values are $(1,0.64,0)$.
Their trapezoidal AUCs are 0.67 and 0.57.
The exact-IG feature-1-first order instead gives insertion AUC 0.57 and deletion AUC 0.67.
The inaccurate order therefore scores better by both conventional criteria.
The transverse term alters coordinate gradients along the diagonal without
altering these finite off-path intervention values. Exact path integration
and perturbation-score optimization are distinct questions.

CPU fp32/fp64 checks at $K=128$ verify trapezoidal coordinates, completeness,
repeatability and responses: budgets 32, 64 and 128 alias; 256 recovers the analytic
order to numerical tolerance. The endpoint-inclusive result is analytic only;
these checks use no RDT data.

\end{document}
